% Generated draft: requires mathematical and source review.
\documentclass[11pt]{article}
\usepackage[margin=25mm]{geometry}
\usepackage{fontspec}
\setmainfont{Latin Modern Roman}
\usepackage{amsmath,amssymb,mathtools,unicode-math}
\setmathfont{Latin Modern Math}
\usepackage{enumitem,needspace,xurl,hyperref}
\setlength{\parindent}{0pt}
\setlength{\parskip}{5pt}
\setlength{\emergencystretch}{3em}
\newcommand{\N}{\mathbb{N}}
\newcommand{\Z}{\mathbb{Z}}
\newcommand{\Q}{\mathbb{Q}}
\newcommand{\R}{\mathbb{R}}
\newcommand{\C}{\mathbb{C}}
\newcommand{\F}{\mathbb{F}}
\newcommand{\A}{\mathbb{A}}
\newcommand{\PP}{\mathbb{P}}
\newcommand{\E}{\mathbb{E}}
\newcommand{\eps}{\varepsilon}
\newcommand{\dd}{\,\mathrm{d}}
\begin{document}
\section*{UnsolvedMath ID: 10400063}
\textbf{Source problem code:} catalogue entry
\subsection{Definitions and mathematical statement}
In the usual Drinfeld setting, an associator consists of a formal series \(\Phi(X,Y)\) in the degree completion of the noncommutative algebra on \(X,Y\), together with a normalization parameter, satisfying the associator pentagon and hexagon identities. Its coefficients lie in a characteristic-zero field \(F\). Group-likeness means
\[
 \Delta\Phi=\Phi\otimes\Phi,\qquad
 \Delta X=X\otimes1+1\otimes X,\qquad
 \Delta Y=Y\otimes1+1\otimes Y,
\]
with the usual completed tensor products and constant term \(1\). The exact parameter, additional normalizations and identity conventions must be taken from the source, which does not give them in this extract.

Problem 3.13 asks for a combinatorial direct presentation of an associator in all degrees, and in particular one with rational coefficients. Thus the presentation must specify a single series
\[
 \Phi=1+\sum_{r\ge1}\Phi_r,\qquad
 \Phi_r=\sum_{|w|=r}a_w w,
\]
whose degree truncations are mutually compatible and satisfy the associator identities to every order. The rational-coefficient version requires \(a_w\in\Q\) for every word \(w\). A construction of unrelated finite-degree approximations, or merely an existence assertion, does not by itself supply the requested direct all-degree presentation. What counts as ``combinatorial'' and ``direct'' is not specified and must not be replaced by an invented algorithmic constraint.

The supplied partial-progress summary claims that the KZ associator has rational coefficients, names combinatorial rational constructions by Bar-Natan and Le--Murakami, and describes the Alekseev--Torossian (2008) associator as combinatorial. These are source-status claims requiring review. In particular, the conventional analytic KZ associator has coefficients expressed through multiple-zeta values with normalization factors; it must be distinguished from the existence of rational associators. The blanket rational-KZ assertion is not adopted as a definition or verified fact.
\subsection{Short English statement}
Give a direct combinatorial construction of one associator valid in every degree, especially with all coefficients rational. The supplied record reports partial progress but does not define directness or give complete construction references; its rational-KZ claim also needs correction or qualification from the source literature.
\subsection{Sources}
\begin{itemize}
\item Local UnsolvedMath registry, permanent ID \(10400063\), exact source code \texttt{catalogue entry}, Problem 3.13. No individual proposer or proposed year is supplied; the title is truncated.
\item Ohtsuki, \emph{Problems on invariants of knots and 3-manifolds} (2002), Problem 3.13, PDF page 61: \url{https://msp.org/gtm/2002/04/gtm-2002-04-024s.pdf}. Printed bold headings distinguish 231 numbered items from surrounding references.
\item Extraction pdf-font-structure, accessed 2026-07-29. Source evidence NEEDS\_REVIEW presents an open problem and calls for a current primary-literature audit. Registry status partially\_solved and classification PARTIAL-PROGRESS name Drinfeld/KZ, Bar-Natan, Le--Murakami and Alekseev--Torossian (2008), without complete bibliographic entries or a precise description of which part is resolved.
\item Dataset attribution: Ulam AI / UnsolvedMath Contributors, CC BY 4.0. Source rights NEEDS\_REVIEW: publicly accessible author/publisher copy, redistribution terms requiring review.
\end{itemize}
\end{document}
